\lecture{6.8}{Reporting a Super Learner Analysis}
\label{lec:6-8}

\subsection{Number of Folds}

A report should state the number of cross-validation folds used to train and evaluate
the Super Learner. Common choices are $V=5$ and $V=10$. If nested cross-validation
was used, report both the outer folds and the inner folds.

For example, one might write:
\begin{quote}
The Super Learner was trained using ten-fold cross-validation to choose ensemble
weights. Predictive performance of the full Super Learner procedure was evaluated
using nested ten-fold cross-validation.
\end{quote}

\subsection{Candidate Algorithms}

The candidate library is part of the estimator and must be reported. Include the
algorithm names, important tuning parameters, and whether any variable screening
procedures were paired with each learner.

\[
\begin{array}{lll}
\toprule
\text{Learner} & \text{Description} & \text{Screener}\\
\midrule
\texttt{SL.glm} & \text{main-terms GLM} & \text{All variables}\\
\texttt{SL.mean} & \text{intercept-only model} & \text{All variables}\\
\texttt{SL.ranger} & \text{random forest} & \text{All variables}\\
\texttt{SL.glm} & \text{main-terms GLM} & \texttt{screen.corP}\\
\bottomrule
\end{array}
\]

Large libraries are often best placed in an appendix table.

\subsection{Ensemble Construction}

State whether the analysis used a discrete Super Learner or an ensemble Super
Learner. For ensemble Super Learner, describe the metalearner. For continuous
outcomes, a convex least-squares ensemble is common. For binary outcomes or
propensity scores, a logistic-linear ensemble minimizing negative log-likelihood is
often appropriate.

\subsection{Loss Functions}

The loss function determines what the Super Learner is optimizing. Reports should
state the loss separately for each nuisance function:
\begin{itemize}
\item outcome regression with continuous $Y$: mean squared-error loss;
\item outcome regression with binary $Y$: Bernoulli negative log-likelihood or
another valid binary-outcome loss;
\item propensity score: Bernoulli negative log-likelihood for treatment prediction.
\end{itemize}

\subsection{Cross-Validated Risk Summaries}

Risk summaries should include the Super Learner, the discrete Super Learner, and the
candidate algorithms. A clear table includes the risk estimate, an uncertainty summary
when available, and the ensemble weight.

\[
\begin{array}{lccc}
\toprule
\text{Learner} & \text{CV risk} & \text{Weight} & \text{Comment}\\
\midrule
\texttt{SL.glm\_All} & 1.03 & 0.40 & \text{stable benchmark}\\
\texttt{SL.mean\_All} & 2.72 & 0.00 & \text{intercept only}\\
\texttt{SL.quadglm\_All} & 1.05 & 0.60 & \text{quadratic terms}\\
\text{Super Learner} & 1.01 & -- & \text{ensemble}\\
\bottomrule
\end{array}
\]

\subsection{Sensitivity Analyses}

Super Learner results depend on the candidate library, loss function, and
cross-validation design. Sensitivity analyses can examine whether conclusions change
when:
\begin{itemize}
\item additional flexible learners are added;
\item unstable or extremely slow learners are removed;
\item the number of folds is changed;
\item required confounders are forced into every candidate;
\item alternative reasonable losses are used;
\item propensity score learners producing extreme weights are excluded or modified.
\end{itemize}

\subsection{Methods Write-Up}

A concise methods paragraph for an outcome regression might read:

\begin{quote}
The outcome regression was estimated using Super Learner. Candidate algorithms
included generalized linear models, penalized regression, generalized additive models,
random forests, and multivariate adaptive regression splines, with selected learners
paired with pre-specified screening rules. Ensemble weights were chosen to minimize
ten-fold cross-validated mean squared-error loss using a convex least-squares
metalearner.
\end{quote}

For a propensity score:

\begin{quote}
The propensity score was estimated using Super Learner with treatment as the outcome
and baseline covariates as predictors. Ensemble weights were chosen to minimize
ten-fold cross-validated Bernoulli negative log-likelihood using a logistic-linear
metalearner.
\end{quote}

\subsection{Results Write-Up}

A results paragraph should summarize which learners contributed meaningfully and how
the ensemble compared with individual candidates:

\begin{quote}
For the outcome regression, the Super Learner placed nonzero weight on five
candidate algorithms, with the largest weights assigned to a random forest and a
generalized additive model. Nested cross-validation suggested that the ensemble had
lower estimated risk than any single candidate. For the propensity score, the ensemble
performed similarly to the best individual learner, and predicted treatment
probabilities showed no severe near-zero or near-one instability.
\end{quote}

Avoid interpreting weights as evidence that a scientific model is true. They describe
predictive contributions under the chosen loss.

\subsection{Appendix Tables}

Appendix tables should make the analysis reproducible. At minimum, include:
\begin{itemize}
\item wrapper function names;
\item plain-language descriptions of each algorithm;
\item tuning parameters;
\item screening rules;
\item cross-validated risks;
\item Super Learner weights;
\item software versions when important.
\end{itemize}

\begin{assumptionbox}
Reporting Super Learner details does not replace reporting causal identification
assumptions. The final causal analysis should separately state the target causal
parameter, consistency, exchangeability, positivity, the adjustment set, and the
statistical estimator.
\end{assumptionbox}

\begin{estimationbox}
A complete Super Learner report lets another analyst reconstruct the nuisance
estimation procedure: folds, library, screeners, losses, metalearner, risk summaries,
weights, and sensitivity analyses.
\end{estimationbox}
